\begin{proof}[Proof of \cref{obj:prop:identification}]
\begin{enumerate}
\item \emph{Completion.}
Fix \(\mathbb P\in\mathcal D_{d,\epsilon}^{\mathrm{obs}}\), and use its cell quantities from \cref{obj:def:observed-class}. Define the Bernoulli probability weights by
\[
\operatorname{Bern}(\vartheta;y)
=
\begin{cases}
1-\vartheta,&y=0,\\
\vartheta,&y=1,
\end{cases}
\qquad
\vartheta\in[0,1],\quad y\in\{0,1\}.
\]
For \(x\in[d]\) and \(a,y,y_0,y_1\in\{0,1\}\), define the full-data atom probabilities by
\[
\begin{aligned}
&P(X=x,A=a,Y=y,Y(0)=y_0,Y(1)=y_1)\\
&\qquad=
\begin{cases}
q_{ay,x}\operatorname{Bern}(\mu_{1-a,x};y_{1-a}),&y=y_a,\\
0,&y\ne y_a,
\end{cases}
\end{aligned}
\]
where
\[
y_a=
\begin{cases}
y_0,&a=0,\\
y_1,&a=1,
\end{cases}
\qquad
y_{1-a}=
\begin{cases}
y_1,&a=0,\\
y_0,&a=1.
\end{cases}
\]
The bounds \(q_{ay,x}\ge0\) and \(\mu_{ax}\in[0,1]\) from \cref{obj:def:observed-class} make these masses nonnegative. Summing first over the potential-outcome coordinates gives
\[
\begin{aligned}
&\sum_{x\in[d]}\sum_{a,y,y_0,y_1\in\{0,1\}}
P(X=x,A=a,Y=y,Y(0)=y_0,Y(1)=y_1)\\
&\quad=
\sum_{x\in[d]}\sum_{a,y\in\{0,1\}}
q_{ay,x}
\bigl\{
\operatorname{Bern}(\mu_{1-a,x};0)
+\operatorname{Bern}(\mu_{1-a,x};1)
\bigr\}\\
&\quad=\sum_{x\in[d]}\sum_{a,y\in\{0,1\}}q_{ay,x}=1.
\end{aligned}
\]
Thus the construction defines a probability law. Its atoms with \(y\ne y_a\) have zero mass, so
\[
P\{Y=Y(A)\}=1,
\]
which is \cref{obj:ass:consistency}.
% lean: CausalSmith.Stat.OptvalueVanishingoverlapRate.completionFullMass_nonneg_sum
% lean: CausalSmith.Stat.OptvalueVanishingoverlapRate.bernoulliCompletion_consistency

Summing over the observed outcome, and then using the observed-atom factorization supplied by \cref{obj:def:observed-class}, yields
\[
\begin{aligned}
&P(X=x,A=a,Y(0)=y_0,Y(1)=y_1)\\
&\quad=q_{a\,y_a,x}\operatorname{Bern}(\mu_{1-a,x};y_{1-a})\\
&\quad=p_x\operatorname{Bern}(\pi_x;a)
\operatorname{Bern}(\mu_{0x};y_0)
\operatorname{Bern}(\mu_{1x};y_1).
\end{aligned}
\]
Consequently, summing these factors over the remaining binary coordinates gives, for \(a,a',y\in\{0,1\}\),
\[
\begin{aligned}
P(X=x)&=p_x,\\
P(X=x,A=a')&=p_x\operatorname{Bern}(\pi_x;a'),\\
P(X=x,Y(a)=y)&=p_x\operatorname{Bern}(\mu_{ax};y),\\
P(X=x,A=a',Y(a)=y)
&=p_x\operatorname{Bern}(\pi_x;a')
\operatorname{Bern}(\mu_{ax};y).
\end{aligned}
\]
Multiplication therefore gives
\[
\begin{aligned}
&P(X=x,A=a',Y(a)=y)\,P(X=x)\\
&\qquad=P(X=x,A=a')\,P(X=x,Y(a)=y).
\end{aligned}
\]
On positive-mass cells, division by \(p_x^2\) proves \cref{obj:ass:conditional-exchangeability}; on null cells both products are zero.
% lean: CausalSmith.Stat.OptvalueVanishingoverlapRate.bernoulliCompletion_poAtom_factorization
% lean: CausalSmith.Stat.OptvalueVanishingoverlapRate.bernoulliCompletion_exchangeability

For each observed atom, the construction also gives
\[
\begin{aligned}
P(X=x,A=a,Y=y)
&=q_{ay,x}
\bigl\{
\operatorname{Bern}(\mu_{1-a,x};0)
+\operatorname{Bern}(\mu_{1-a,x};1)
\bigr\}\\
&=q_{ay,x}.
\end{aligned}
\]
Hence
\[
\operatorname{Obs}(P)=\mathbb P.
\]
The observed marginal inherits the overlap restriction of \(\mathbb P\). Together with consistency and conditional exchangeability, this puts \(P\) in the class of \cref{obj:def:causal-class}, proving the completion assertion.
% lean: CausalSmith.Stat.OptvalueVanishingoverlapRate.bernoulliCompletion_observedMarginal

\item \emph{Value identification.}
Fix \(P\in\mathcal P_{d,\epsilon}\), and use \(q_{ay,x},p_x,\pi_x,\mu_{ax}\) for the quantities of its observed marginal defined in \cref{obj:def:observed-class}. Marginalization gives
\[
P(X=x)
=\sum_{a,y\in\{0,1\}}
\operatorname{Obs}(P)(X=x,A=a,Y=y)
=p_x.
\]
We compare the two value sums cell by cell. If \(p_x=0\), their cell contributions are both zero:
\[
p_x\max_{a\in\{0,1\}}\mathbb E_P[Y(a)\mid X=x]
=
p_x\max_{a\in\{0,1\}}\mu_{ax}
=0.
\]
Here the conditional potential-outcome means use the stated zero convention.
% lean: CausalSmith.Stat.OptvalueVanishingoverlapRate.observedMarginal_cellMass

Suppose now that \(p_x>0\). The lower overlap bound in \cref{obj:ass:shrinking-overlap} gives
\[
q_{10,x}+q_{11,x}=p_x\pi_x\ge\epsilon p_x>0.
\]
The decomposition of the cell mass and the upper overlap bound give
\[
\begin{aligned}
p_x&=(q_{00,x}+q_{01,x})+(q_{10,x}+q_{11,x}),\\
q_{10,x}+q_{11,x}&\le(1-\epsilon)p_x,
\end{aligned}
\]
and hence
\[
q_{00,x}+q_{01,x}
=p_x-(q_{10,x}+q_{11,x})
\ge\epsilon p_x>0.
\]
Thus both arm masses are positive.

For either \(a\in\{0,1\}\), \cref{obj:ass:conditional-exchangeability} supplies
\[
P(X=x,A=a,Y(a)=1)\,p_x
=
P(X=x,A=a)\,P(X=x,Y(a)=1).
\]
Consistency, as specified in \cref{obj:ass:consistency}, identifies the selected atoms:
\[
P(X=x,A=a,Y(a)=y)=q_{ay,x},
\qquad y\in\{0,1\}.
\]
Summing this identity over \(y\), and substituting its outcome-one instance into the exchangeability identity, gives
\[
P(X=x,A=a)=q_{a0,x}+q_{a1,x},
\qquad
q_{a1,x}p_x
=(q_{a0,x}+q_{a1,x})P(X=x,Y(a)=1).
\]
Division by the positive cell and arm masses now yields
\[
\mathbb E_P[Y(a)\mid X=x]
=
\frac{P(X=x,Y(a)=1)}{p_x}
=
\frac{q_{a1,x}}{q_{a0,x}+q_{a1,x}}
=\mu_{ax}.
\]
Combining the positive- and zero-mass cases and summing over \(x\) proves
\[
\begin{aligned}
V^\star(P)
&=\sum_{x\in[d]}P(X=x)
\max_{a\in\{0,1\}}\mathbb E_P[Y(a)\mid X=x]\\
&=\sum_{x\in[d]}p_x\max_{a\in\{0,1\}}\mu_{ax}
=\Psi\{\operatorname{Obs}(P)\},
\end{aligned}
\]
where the last equality is the definition in \cref{obj:def:observed-value}.
% lean: CausalSmith.Stat.OptvalueVanishingoverlapRate.poRegression_eq_outcomeMean_of_pos
% lean: CausalSmith.Stat.OptvalueVanishingoverlapRate.identification

\item \emph{The armwise extension on observed cells.}
Fix \(\mathbb P\in\mathcal D_{d,\epsilon}^{\mathrm{obs}}\) and \(x\in[d]\). The cell vector and the quantities in \cref{obj:def:armwise-extension} satisfy
\[
s(q_x)=p_x,\qquad
s_a(q_x)=q_{a0,x}+q_{a1,x},\qquad
(q_x)_{a1}=q_{a1,x},
\qquad a\in\{0,1\}.
\]
Every coordinate is nonnegative, so \(p_x\ge0\). If \(p_x=0\), then
\[
0\le q_{ay,x}\le p_x=0
\qquad(a,y\in\{0,1\}),
\]
and therefore \(q_x=0\). The zero branch of \cref{obj:def:armwise-extension} gives
\[
F_\epsilon(q_x)=0
=p_x\max_{a\in\{0,1\}}\mu_{ax}.
\]

Suppose instead that \(p_x>0\). The lower overlap bound gives
\[
\epsilon p_x\le q_{10,x}+q_{11,x}=s_1(q_x).
\]
For the other arm, the upper overlap bound and the cell-mass decomposition give
\[
s_1(q_x)\le(1-\epsilon)p_x,
\qquad
s_0(q_x)=p_x-s_1(q_x)\ge\epsilon p_x.
\]
Thus, for each \(a\in\{0,1\}\),
\[
s_a(q_x)\ge\epsilon p_x>0,
\qquad
D_a(q_x)=\max\{s_a(q_x),\epsilon p_x\}=s_a(q_x).
\]
Each armwise term consequently reduces to
\[
\frac{s(q_x)(q_x)_{a1}}{D_a(q_x)}
=
\frac{p_xq_{a1,x}}{q_{a0,x}+q_{a1,x}}
=
p_x\mu_{ax}.
\]
Since \(s(q_x)=p_x>0\), the vector \(q_x\) is nonzero. Using the nonzero branch of \cref{obj:def:armwise-extension} and factoring the nonnegative \(p_x\) out of the maximum gives
\[
F_\epsilon(q_x)
=
\max_{a\in\{0,1\}}p_x\mu_{ax}
=
p_x\max_{a\in\{0,1\}}\mu_{ax}.
\]
This proves the displayed armwise identity, including its null-cell case.
% lean: CausalSmith.Stat.OptvalueVanishingoverlapRate.armwiseExtension_cellVector
\end{enumerate}
\end{proof}
